\documentclass[fleqn]{report}

\usepackage[serbian]{babel}
\usepackage[a4paper]{geometry}
\usepackage{amsmath}
\usepackage{csquotes}
\usepackage{amssymb}
\usepackage{siunitx}
\sisetup{ per-mode = symbol }

\usepackage{cancel}
\usepackage{empheq}
\usepackage{pgfplots}
\pgfplotsset{compat=1.18}
\usepgfplotslibrary{external}
\usetikzlibrary{arrows.meta, angles, quotes, patterns, patterns.meta}
\tikzexternalize
\tikzset{
    >=Stealth,
}

\title{
    Fizička elektronika gasova i plazme \\
    \large rešeni zadaci za prvi kolokvijum
}
\author{}
\date{Školska 2026/27. godina}

\begin{document}

\maketitle

\begin{enumerate}
    \item Izvesti izraz za Debye-vu dužinu određivanjem potencijala tačkastog naelektrisanja $Q$ u plazmi.
          Elektroni se raspodeljuju po Bolcmanovoj raspodeli, dok se može smatrati da je koncentracija jona homogena.
          Smatrati da važi gasna aproksimacija.
          Pokazati da je broj elektrona unutar Debye-ve sfere veoma veliki.

          \underline{Rešenje}:

          Koncentracija elektrona je po Bolcmanovoj raspodeli data sa:
          \begin{align*}
               & n_e = n_0 \exp\left(-\frac{U}{k_B T}\right), \quad U = -eV \\
               & n_e = n_0 \exp\left(\frac{eV}{k_B T}\right)
          \end{align*}
          Kada naelektrisanje nije prisutno, $V = 0$, pa je $n_e = n_0$, a samim tim i $n_i = n_0$ (zbog kvazineutralnosti plazme $n_e = n_i$).
          Pretpostavljamo da su joni dovoljno inertni tako da i u prisustvu naelektrisanja $Q$ ostanu jednako raspodeljeni, tj. da ostane $n_i = n_0$.
          Sa poznatim izrazima za raspodele naelektrisanja možemo tražiti potencijal iz Puasonovog zakona:
          \begin{align*}
               & \nabla \cdot \vec E = \frac{\rho}{\varepsilon_0} \implies \nabla^2 V = -\frac{\rho}{\varepsilon_0}                                                                                                                   \\
               & \begin{cases}
                     \nabla^2 = \frac{\partial^2}{\partial x^2} + \frac{\partial^2}{\partial y^2} + \frac{\partial^2}{\partial z^2} \\
                     \rho = n_i e - n_e e
                 \end{cases}                                                                                                   \\
               & \nabla^2 V = -\frac{e}{\varepsilon_0} (n_i - n_e) = \frac{e}{\varepsilon_0} \left(n_0\exp\left(\frac{eV}{k_B T}\right) - n_0\right) = \frac{n_0 e}{\varepsilon_0} \left(\exp\left(\frac{eV}{k_B T}\right) - 1\right)
          \end{align*}
          U gasnoj aproksimaciji je kinetička energija mnogo veća od potencijalne, tj. $k_B T \gg eV$, a za malo $x$ je $\exp(x) \approx 1 + x$, pa možemo pisati:
          \[ \nabla^2 V \approx \frac{n_0 e}{\varepsilon_0} \left(1 + \frac{eV}{k_B T} - 1\right) = \frac{n_0 e^2}{\varepsilon_0 k_B T} V \]
          Ako uvedemo oznaku $\lambda_D = \sqrt{\frac{\varepsilon_0 k_B T}{n_0 e^2}}$:
          \[ \nabla^2 V = \frac{1}{\lambda_D^2} V \]
          Ako pretpostavimo sferno simetrično rešenje, tako da je $\frac{\partial V}{\partial \theta} = \frac{\partial V}{\partial \phi} \equiv 0$, onda se laplasijan u sfernom koordinatnom sistemu značajno pojednostavljuje i dobijamo:
          \[ \frac{1}{r^2} \frac{\partial}{\partial r} \left(r^2 \frac{\partial V}{\partial r}\right) = \frac{1}{\lambda_D^2} V \]
          Ovu jednačinu rešavamo uvođenjem smene $V(r) = \frac{F(r)}{r}$:
          \begin{align*}
               & \frac{\partial V}{\partial r} = \frac{\frac{dF}{dr} \cdot r - F}{r^2} = \frac{1}{r}\frac{dF}{dr} - \frac{F}{r^2}                                  \\
               & r^2 \frac{\partial V}{\partial r} = r\frac{dF}{dr} - F                                                                                            \\
               & \frac{\partial}{\partial r} \left(r^2 \frac{\partial V}{\partial r}\right) = \cancel{\frac{dF}{dr}} + r\frac{d^2F}{dr^2} - \cancel{\frac{dF}{dr}} \\
               & \frac{1}{r^2}\frac{\partial}{\partial r} \left(r^2 \frac{\partial V}{\partial r}\right) = \frac{1}{r}\frac{d^2F}{dr^2}
          \end{align*}
          Ubacivanjem ovog izraza u gornju jednačinu dobijamo:
          \begin{align*}
               & \frac{1}{\cancel{r}}\frac{d^2F}{dr^2} = \frac{1}{\lambda_D^2} \frac{F(r)}{\cancel{r}}           \\
               & \implies F(r) = A\exp\left(-\frac{r}{\lambda_D}\right) + B\exp\left(+\frac{r}{\lambda_D}\right)
          \end{align*}
          Za veliko rastojanje, očekujemo da uronjeno naelektrisanje nema nikakav efekat, dakle da se potencijal spusti na 0:
          \[ F(r \to \infty) \to 0 \implies B = 0 \]
          Slično, za dovoljno malo rastojanje od tačkastog naelektrisanja, elektroni koji su privučeni ne menjaju značajno oblik polja, pa očekujemo da za malo $r$ bude:
          \begin{align*}
               & V(r \to 0) \to \frac{Q}{4\pi\varepsilon_0 r} \implies F(r \to 0) \to \frac{Q}{4\pi\varepsilon_0} = A \\
               & \implies V = \frac{Q}{4\pi\varepsilon_0 r} \exp\left(-\frac{r}{\lambda_D}\right)
          \end{align*}

          \noindent\hfill
          \begin{tikzpicture}
              \begin{axis}[
                      axis lines=middle,
                      xmin=-0.1, xmax=5,
                      ymin=-0.1, ymax=5,
                      xlabel={$r$},
                      ylabel={$V$},
                      xtick={1},
                      xticklabels={$\lambda_D$},
                      ytick=\empty,
                  ]
                  \addplot[
                      domain=0.2:5,
                      samples=200,
                      thick,
                  ]{exp(-x)/x};
                  \addplot[
                      domain=0.2:4.5,
                      samples=200,
                      thick,
                      dotted,
                  ]{1/x}node[pos=0.5, above right] {$\sim\frac{1}{r}$};;
              \end{axis}
          \end{tikzpicture}
          \hfill\null\\
          Na osnovu grafika se vidi da se na rastojanju većem od $\lambda_D$ uronjeno naelektrisanje praktično ne vidi.

          Ostalo je još da pokažemo da je broj elektrona u Debye-voj sferi veliki, odnoso da pokažemo da važi $n_0 \lambda_D^3 \gg 1$.
          Krenimo od gasne aproksimacije, koja kaže da je kinetička energija čestica mnogo veća od potencijalne:
          \begin{align*}
               & \frac{3}{2}k_B T_e \gg \frac{e^2}{4\pi\varepsilon_0 \overline{d}} \text{ ($\overline{d}$ je srednje rastojanje između čestica)}                 \\
               & k_B T_e \gg \frac{e^2}{\varepsilon_0 \overline{d}}                                                                                              \\
               & \frac{\varepsilon_0 k_B T_e \overline{d}}{e^2} \gg 1                                                                                            \\
               & n_0 \approx \frac{1}{\overline{d}^3} \implies \overline{d} \approx n_0^{-\frac{1}{3}}                                                           \\
               & \frac{\varepsilon_0 k_B T_e}{e^2 n_0^{-\frac{1}{3}}} = \underbrace{\frac{\varepsilon_0 k_B T_e}{n_0 e^2}}_{\lambda_D^2} n_0^{\frac{2}{3}} \gg 1 \\
               & \lambda_D^2 n_0^{\frac{2}{3}} \gg 1 \bigg/ ^{\frac{3}{2}}                                                                                       \\
               & \lambda_D^3 n_0 \gg 1
          \end{align*}
          \hrulefill
    \item Odrediti i skicirati strujno-naponsku karakteristiku Langmuir-ove sonde uronjene u plazmu.
          Objasniti na koji način se na osnovu strujno-naponske karakteristike mogu odrediti temperatura i koncentracija elektrona.

          \underline{Rešenje}:

          \emph{
              Langmuir-ova sonda je instrument za merenje parametara plazme.
              U suštini predstavlja provodni objekat kome možemo podešavati potencijal i uroniti ga u plazmu.
              Elektroni koji padaju na provodnu površ čine struju kroz sondu koju možemo meriti.
              Posmatranjem zavisnosti izmerene struje od potencijala $I(V)$ možemo zaključiti o kojim se vrednostima parametara radi.
          }

          Označimo sa $V$ potencijal sonde, $V_p$ potencijal plazme u okolini sonde, i $V_s$ razliku tih potencijala $V_s = V - V_p$.

          Posmatramo prvo slučaj kada je $V_s$ negativno.
          Tada struju čine elektroni samo elektroni koji imaju dovoljno veliku energiju da preskoče potencijalnu barijeru koja se formira između sonde i plazme:
          \[ \frac{m_e v_z^2}{2} > |eV_s| \implies v_z > \sqrt{\frac{2|eV_s|}{m_e}} \quad \left(\implies p_z > \sqrt{2m_e|eV_s|} = p_{z1}\right) \]
          $v_z$ je brzina elektrona u pravcu normalnom na površ sonde.
          Gustina struje tih elektrona je:
          \[ \vec{J} = ne\left<\vec{v}\right> \]
          pa je ukupna struja:
          \[ I = \int \vec{J} d\vec{S} = \vec{J} \vec{S} = e n_e S \left<v_z\right> \]
          gde je S uronjena površina sonde.
          \begin{align*}
              \left<v_z\right> & = \int v_z f_1(\vec{r}, \vec{p}) d\vec{r} d\vec{p} =                \\
                               & = \int \frac{p_z}{m} f_M(\vec{p}) d\vec{p}                          \\
              f_M(\vec{p})     & = (2 \pi k_B T_e m_e)^{-\frac{3}{2}} e^{-\frac{p^2}{2 m_e k_B T_e}}
          \end{align*}
          \begin{align*}
              I_e & = e n_e S \int_{p_x = -\infty}^{+\infty} \int_{p_y = -\infty}^{+\infty}\int_{p_z = p_{z1}}^{+\infty} \frac{p_z}{m_e} (2 \pi k_B T_e m_e)^{-\frac{3}{2}} e^{-\frac{p_x^2 + p_y^2 + p_z^2}{2m_e k_B T}} dp_x dp_y dp_z =                                              \\
                  & = \frac{e n_e S}{m_e} (2 \pi k_B T_e m_e)^{-\frac{3}{2}} {\underbrace{\left(\int_{-\infty}^{+\infty} e^{-\frac{p_x^2}{2m_e k_B T_e}} dp_x\right)}_{\sqrt{2\pi m_e k_B T_e}\footnotemark}}^2 \cdot \int_{p_{z1}}^{+\infty} p_z e^{-\frac{p_z^2}{2m_e k_B T_e}}dp_z = \\
                  & = \frac{e n_e S}{m_e} (2 \pi k_B T_e m_e)^{-\frac{1}{2}} \int_{p_{z1}}^{+\infty} p_z e^{-\frac{p_z^2}{2m_e k_B T_e}}dp_z = \bigg/ \begin{matrix}
                                                                                                                                                          u = -\frac{p_z^2}{2 m_e k_B T_e} \\
                                                                                                                                                          du = -\frac{\cancel 2 p_z dp_z}{\cancel 2 m_e k_B T_e}
                                                                                                                                                      \end{matrix} \bigg/                                                                         \\
                  & = \frac{e n_e S}{\cancel{m_e}} (2 \pi k_B T_e m_e)^{-\frac{1}{2}} (-\cancel{m_e} k_B T_e) \int_{u_1}^{-\infty} e^u du =                                                                                                                                             \\
                  & = e n_e S \frac{1}{\sqrt{2 \pi k_B T_e m_e}} k_B T_e \left. e^u \right|_{-\infty}^{u_1} =                                                                                                                                                                           \\
                  & = e n_e S \underbrace{\sqrt{\frac{k_B T_e}{2 \pi m_e}}}_{\frac{\overline{v}}{4}\footnotemark} e^{u_1} =                                                                                                                                                             \\
                  & = \underbrace{e n_e S \frac{\overline{v}}{4}}_{I_{es}} e^{-\frac{|eV_s|}{k_B T}} =                                                                                                                                                                                  \\
              I_e & = I_{es} \exp\left(\frac{e(V - V_p)}{k_B T}\right)
          \end{align*}
          \footnotetext{$\int_{-\infty}^{+\infty} e^{-\alpha x^2} dx = \sqrt{\frac{\pi}{\alpha}}$}
          \footnotetext{$\overline{v} = \sqrt{\frac{8 k_B T_e}{\pi m_e}}$}
          Ako se ova strujno naponska karakteristika logaritmuje, dobija se linearna zavisnost:
          \[ \ln I_e = \ln I_{es} + \frac{eV}{k_B T_e} - \frac{e V_p}{k_B T_e} \]
          Posmatranjem nagiba ove karakteristike može se zaključiti temperatura elektrona u plazmi, a posmatranjem slobodnog člana, ako je $V_p$ poznato, može se doći do struje zasićenja $I_{es}$, a odatle i koncentracije elektrona.

          Za pozitivne vrednosti $V_s$ sonda privlači elektrone, tako da će jako brzo potrošiti sve elektrone na raspolaganju u svojoj brzini, pa će struja kroz nju biti ograničena brzinom generisanja ili dovođenja novih elektrona.
          Na grafiku se ovaj efekat vidi kao zasićenje krive za vrednosti $V > V_p$, dok se u okolini $V = V_p$ javlja koleno, pa se iz pozicije kolena može odrediti i $V_p$.

          \noindent\hfill
          \begin{tikzpicture}
              \begin{axis}[
                      axis lines=middle,
                      xmin=-5, xmax=5,
                      ymin=-2, ymax=5,
                      xlabel={$V$},
                      ylabel={$I_e$},
                      xtick={2},
                      xticklabels={$V_p$},
                      ytick={2.216},
                      yticklabels={$I_{es}$},
                  ]
                  \addplot[
                      domain=-5:2,
                      samples=200,
                      thick,
                  ]{0.3*exp(x)};
                  \addplot[
                      domain=2:5,
                      samples=200,
                      thick,
                  ]{0.3*sqrt(x-1.99) + 2.1867};
                  \draw[dotted] (2,0)--(2,2.216);
                  \draw[dotted] (0,2.216)--(2,2.216);
              \end{axis}
          \end{tikzpicture}
          \hfill\null\\

          \hrulefill
    \item Posmatrajmo dve beskonačne ravni na rastojanju $2d$ i potencijalu $V = 0$.
          Prostor je popunjen naelektrisanjima $q$ koncentracije $n$.
          Izračunati potencijal između dve ravni.
          Pokazati da je za $d \gg \lambda_D$ energija potrebna za transport čestica od jedne ravni do sredine veća ili jednaka od srednje kinetičke energije $\left<E_k\right> = \frac{k_B T}{2}$

          \underline{Rešenje}:

          \noindent\hfill
          \begin{tikzpicture}
              \begin{axis}[
                      axis lines=middle,
                      clip mode=individual,
                      xlabel={$x$},
                      xmin=-3, xmax=3,
                      ymin=-3, ymax=3,
                      xtick={-2, 0, 2},
                      xticklabels={$-d$, $0$, $d$},
                      ytick=\empty,
                      axis y line=none
                  ]
                  \draw[thick] (axis cs:-2,-2) -- (axis cs:-2,2);
                  \draw[thick] (axis cs:2,-2) -- (axis cs:2,2);
                  \fill[pattern=dots] (axis cs: -2, -2) rectangle (axis cs: 2, 2);
              \end{axis}
          \end{tikzpicture}
          \hfill\null\\
          Zavisnost potencijala tražimo iz Puasonovog zakona:
          \begin{align*}
               & \nabla^2 V = -\frac{\rho}{\varepsilon_0} \implies \frac{d^2 V}{dx^2} = -\frac{qn}{\varepsilon_0} \\
               & \frac{dV}{dx} = -\frac{qn}{\varepsilon_0}x + C_1                                                 \\
               & V(x) = - \frac{qn}{2\varepsilon_0} x^2 + C_1 x + C_2                                             \\
               & V(x = -d) = V(x = +d) = 0                                                                        \\
               & \begin{rcases}
                     -\frac{qn}{2\varepsilon_0}d^2 - C_1 d + C_2 = 0 \\
                     -\frac{qn}{2\varepsilon_0}d^2 + C_1 d + C_2 = 0
                 \end{rcases} \implies C_1 = 0, C_2 = \frac{qnd^2}{2\varepsilon_0}                                \\
               & \implies V(x) = \frac{qn}{2\varepsilon_0}(d^2 - x^2)
          \end{align*}

          Energija koja je potrebna za transport je poporcionalna razlici potencijala:
          \[ E = q\Delta V = q(V(0) - V(d)) = \frac{q^2nd^2}{2\varepsilon_0} \]
          Ako je $d \gg \lambda_D$ onda je i:
          \[ E \gg \frac{q^2n\lambda_D^2}{2\varepsilon_0} = \frac{\cancel{q^2n}}{2\cancel{\varepsilon_0}} \frac{\cancel{\varepsilon_0} k_B T}{\cancel{q^2n}} = \frac{k_B T}{2} \]

          \hrulefill
    \item U TFTR u Prinstonu je zagrejana plazma pomoću ubacivanja neutralnih deuterijumovih atoma $^2_1D$ energije $200\si{\kilo\electronvolt}$.
          Posle ulaska u magnetno polje konvertuju se u $^2_1D^+$ jone.
          Ovi joni mogu da se konfiniraju samo ako je Larmorov poluprečnik $R_L < 0.6\si{\meter}$.
          Izračunati $R_L$ za magnetsku indukciju $B=5\si{\tesla}$ i proveriti da li dolazi do konfinacije jona.

          \underline{Rešenje}:

          Larmorov poluprečnik je dat sa:
          \[ R_L = \frac{v_\perp}{\omega_c} \]
          gde je $\omega_c = \frac{eB}{m}$ ciklotronska frekvencija.
          Proverimo konfinaciju za čestice u najnepovoljnijoj situaciji --- one kod kojih je kompletan vektor brzine normalan na pravac magnetnog polja (imaju najveće $v_\perp$ a samim tim i $R_L$):
          \[ v_\perp = v = \sqrt{\frac{2E_k}{m}} \]
          Dobijamo da je Larmorov poluprečnik u najgorem slučaju:
          \[ R_L = \frac{v_\perp}{\omega_c} = \frac{mv}{eB} = \frac{m}{eB}\sqrt{\frac{2E_k}{m}} = 1.8\si{\centi\meter} \]
          što je mnogo manje od zahtevanih $0.6\si{\meter}$, pa zaključujemo da će do konfinacije jona doći.

          \hrulefill
    \item Neneutralni snop elektrona ima gustinu $n_e = 10^{14}\si{\meter^{-3}}$ i poluprečnik $a = 1\si{\centi\meter}$.
          Snop se kreće duž magnetskog polja indukcije $B = 2\si{\tesla}$.
          Izračunati intenzitet i smer drifta.

          \underline{Rešenje}:

          \noindent\hfill
          \begin{tikzpicture}
              \draw[->] (-4, 0) -- (4, 0) node[right] {$\vec B = B_0 \vec{e}_z$};
              \draw (-3, -2) -- (3, -2);
              \draw (-3, 2) -- (3, 2);
              \draw[->] (3.1, 0) -- (3.1, 2) node[right, pos=0.5] {$a$};
              \draw (-0.5, -1) rectangle (0.5, 1);
              \draw[<->] (-0.5, -1.1) -- (0.5, -1.1) node[below, pos=0.5] {$\Delta z$};
              \draw[->] (0.6, 0) -- (0.6, 1) node[right, pos=0.5] {$r$};
              \draw[->] (0, 0.5) -- (0, 1.5) node[right, pos=0.8] {$\vec{E}$};
              \fill[pattern={Dots[distance=6pt, radius=0.3pt]}] (-3,-2) rectangle (3, 2);
          \end{tikzpicture}
          \hfill\null\\

          Jedini drift koji je prisutan je $E \times B$ drift, gde je brzina drifta data sa:
          \[ \vec v = \frac{\vec E \times \vec B}{B^2} \]
          Vektor $\vec B$ je već poznat, pa je potrebno samo odrediti izraz za vektor električnog polja $\vec E$.
          Pošto je raspodela naelektrisanja data, to možemo naći pomoću Gausovog zakona (posmatramo cilindričnu površ unutar snopa visine $\Delta z$ i poluprečnika $r$):
          \begin{align*}
               & \int_S \vec{E}\cdot d\vec{S} = \frac{1}{\varepsilon_0} \int_V \rho dV                                                          \\
               & E \cdot 2 \cancel{\pi} \cancel r \cancel{\Delta z} = -\frac{e n_e r^{\cancel 2} \cancel{\pi} \cancel{\Delta z}}{\varepsilon_0} \\
               & E = - \frac{en_e}{2\varepsilon_e}r                                                                                             \\
               & \vec{E} = -\frac{e n_e}{2\varepsilon_0} r \vec{e}_r
          \end{align*}
          Sada možemo odrediti i brzinu drifta:
          \[ \vec{v} = \frac{\vec{E}\times\vec{B}}{B^2} = \frac{|\vec{E}|}{B_0} \underbrace{(\vec{e}_r \times \vec{e}_z)}_{-\vec{e}_\phi} = \frac{e n_e r}{2\varepsilon_0 B_0} \vec{e}_\phi \]
          Maksimalna vrednost brzine se dostiže na ivicama snopa:
          \[ v_{max} = v(r = a) = \frac{e n_e a}{2 \varepsilon_0 B_0} = 4517\si{\kilo\meter\per\second} \]
          Možemo da uočimo da elektroni u snopu driftuju \enquote{spiralno}, pa se radi o takozvanom Brillouin-ovom driftu.

          \hrulefill
    \item Data je cilindrična radijalno simetrična plazmena struktura koja se nalazi u homogenom i stacionarnom magnetskom polju $\vec{B} = B \vec{e}_z$.
          Neka je gustina elektrona $n(r) = n_0 \exp\left(-\frac{r^2}{r_0^2}\right)$.
          Pokazati da su brzine drifta usled ukrštenog električnog i magnetskog polja $v_E$ i dijamagnetski drift $v_D$ jednaki i suprotnog smera.

          \underline{Rešenje}:

          Odredimo prvo $v_E$.
          Za razliku od prethodnog zadatka, ovde se ne podrazumeva da su prisutni isključivo elektroni, već u ovoj plazmi mogu biti prisutni i joni, a njihova raspodela nam nije poznata.
          Zato električno polje moramo naći na drugačiji način --- prvo ćemo odrediti potencijal.
          Ako pretpostavimo da su elektroni u plazmi u termodinamičkoj ravnoteži, možemo pisati da zadovoljavaju Bolcmanovu raspodelu:
          \begin{align*}
               & n = n_0 \exp\left(-\frac{r^2}{r_0^2}\right) = n_0 \exp(-\frac{U}{k_B T}) = n_0 \exp(\frac{eV}{k_B T})                                              \\
               & -\frac{r^2}{r_0^2} = \frac{eV}{k_B T}                                                                                                              \\
               & V = -\frac{k_B T}{e}\frac{r^2}{r_0^2}                                                                                                              \\
               & \vec E = -\nabla V = -\frac{\partial V}{\partial r} \vec{e}_r = \frac{k_B T}{e} \frac{2r}{r_0^2} \vec{e}_r                                         \\
               & \vec{v}_E = \frac{\vec{E}\times\vec{B}}{B^2} = \frac{k_B T 2r}{e B_0 r_0^2} \vec{e}_r \times \vec{e}_z = -\frac{2k_B T r}{e B_0 r_0^2}\vec{e}_\phi \\
          \end{align*}
          Sada ćemo odrediti $v_D$, pri čemu ćemo se ponovo pozvati na pretpostavku da je plazma u TDR i reći da je temperatura konstantna:
          \begin{align*}
              \vec{v}_D & = \frac{\vec{B} \times \nabla p}{qnB_0^2} =                                              \\
                        & = \frac{\cancel{B_0} \vec{e}_z \times \nabla (n k_B T)}{-e n B_0^{\cancel 2}} =          \\
                        & = -\frac{k_B T}{e n B_0} \vec{e}_z \times \nabla n =                                     \\
                        & = -\frac{k_B T}{e n B_0} \vec{e}_z \times \left(\frac{dn}{dr} \vec{e}_r\right) =         \\
                        & = -\frac{k_B T}{e \cancel n B_0} \cancel n \left(-\frac{2r}{r_0^2}\right) \vec{e}_\phi = \\
                        & = \frac{2 k_B T r}{e B_0 r_0^2} \vec{e}_\phi = -\vec{v}_E
          \end{align*}

          \hrulefill
    \item Uniformna plazma se nalazi u toroidalnoj komori kvadratnog poprečnog preseka.
          Magnetsko polje je generisano usled struje $I$ duž ose simetrije.
          Dimenzije su $a = 1\si{\centi\meter}$ i $R = 10\si{\centi\meter}$.
          Raspodela u plazmi je Maksvelova $k_B T = 100\si{\electronvolt}$ i ima koncentraciju $n = 10^{19}\si{\meter^{-3}}$.
          Nema električnog polja.
          Izračunati brzinu akumulacije (fluks) naelektrisanih čestica na gornju površinu komore usled $v_{\nabla B}$ i $v_R$ drifta.
          Magnetsko polje u centru komore je jednako $1\si{\tesla}$.

          \underline{Rešenje}:

          Ako je $\partial q$ ukupno naelektrisanje koje padne na gornju površ u vremenskom intervalu $\partial t$, fluks naelektrisanja je:
          \[ \frac{\partial q}{\partial t} = \int \vec{J}d\vec{S} = envS \]
          gde je $S = 2R\pi a$ i $v$ ukupna brzina drifta (usled gradijenta magnetskog polja i usled zakrivljenosti njegovih linija).
          Prvo, potrebni izraz za magnetsko polje dobijamo iz Amperovog zakona:
          \begin{align*}
               & \oint \vec{B} d\vec{l} = \mu_0 I                                                                  \\
               & B \cdot 2r\pi = \mu_0 I                                                                           \\
               & B = \frac{\mu_0 I}{2\pi r} \vec{e}_\phi                                                           \\
               & \nabla B = \frac{dB}{dr} \vec{e}_r = -\frac{\mu_0 I}{2\pi r^2} \vec{e}_r = -\frac{B}{r} \vec{e}_r
          \end{align*}
          Sada odredimo drift usled gradijenta magnetskog polja:
          \begin{align*}
              \vec{v}_{\nabla B} & = \frac{R_L v_\perp}{2B^2} \vec{B} \times \nabla B; \quad R_L = \frac{v_\perp}{\omega_c} = \frac{m v_\perp}{qB}        \\
              \vec{v}_{\nabla B} & = \frac{m v_\perp}{qB} \frac{v_\perp}{2\cancel{B^2}} (\cancel{B} \vec{e}_\phi)\times(-\frac{\cancel{B}}{r}\vec{e}_r) = \\
                                 & = \frac{\frac{m v_\perp^2}{2}}{qBr} (-\vec{e}_\phi \times \vec{e}_r) =                                                 \\
                                 & = \frac{E_\perp}{qBr} \vec{e}_z =                                                                                      \\
                                 & = \frac{k_B T}{qBR} \vec{e}_z \footnotemark
          \end{align*}
          \footnotetext{$\left<E\right> = \frac{3}{2} k_B T = \frac{p^2}{2m} = \frac{p_x^2}{2m} + \frac{p_y^2}{2m} + \frac{p_z^2}{2m} \implies \frac{p_x^2}{2m} = \frac{p_y^2}{2m} = \frac{p_z^2}{2m} = \frac{k_B T}{2} \implies E_\parallel = \frac{k_B T}{2}, E_\perp = k_B T$}
          Odredimo i brzinu drifta usled zakrivljenosti linija magnetskog polja.
          Taj drift možemo da posmatramo kao uopštenje $\vec{E} \times \vec{B}$ drifta, u kom umesto elektrostatiče sile na česticu deluje neka proizvoljna, koja može biti i centrifugalna kada centar vođenja skreće po zakrivljenim linijama polja:
          \begin{align*}
               & \vec{v}_{EB} = \frac{\vec{E}\times\vec{B}}{B^2} \rightarrow \vec{v} = \frac{1}{q}\frac{\vec{F}\times\vec{B}}{B^2}                                                                                  \\
               & \vec{F} = \vec{F}_{cf} = -m\vec{a}_n = -m(-\frac{v_\parallel^2}{R}\vec{e}_r) = \frac{mv_\parallel^2}{R}\vec{e}_r = \frac{2E_\parallel}{R}\vec{e}_r = \frac{k_B T}{R} \vec{e}_r                     \\
               & \vec{v}_R = \frac{1}{q} \frac{\left(\frac{k_B T}{R}\vec{e}_r\right) \times(\cancel{B}\vec{e}_\phi)}{B^{\cancel 2}} = \frac{k_B T}{qBR} \vec{e}_r \times \vec{e}_\phi = \frac{k_B T}{qBR} \vec{e}_z
          \end{align*}
          Ukupna brzina drifta je:
          \begin{align*}
               & \vec{v}_{tot} = \vec{v}_{\nabla B} + \vec{v}_R = \frac{2 k_B T}{qBR} \vec{e}_z                                                                                         \\
               & \frac{\partial q}{\partial t} = \cancel{e} n \frac{2 k_B T}{\cancel{q} \cancel{R} B} \cdot 2 \cancel{R} \pi a = \frac{4 k_B T n \pi a}{B} = 20\si{\coulomb\per\second}
          \end{align*}
          \hrulefill
    \item Proton energije $1\si{\kilo\electronvolt}$ i $v_\parallel = 0$ u homogenom magnetnom polju $B_0 = 0.1\si{\tesla}$ se ubrzava tako što se magnetsko polje lagano povećava na $1\si{\tesla}$.
          Zatim se proton elastično sudara sa teškom česticom i promeni pravac kretanja tako da je $v_\parallel' = v_\perp'$.
          Zatim se polje lagano smanji na $B_0 = 0.1\si{\tesla}$.
          Kolika je sada energija protona?

          \underline{Rešenje}:

          U vremenski promenljivom magnetnom polju naelektrisana čestica zadržava svoj početni magnetni moment:
          \begin{align*}
              \mu = \frac{E_\perp}{B} = const                                                     \\
              \frac{E_{\perp 0}}{B_0} = \frac{E_\perp}{B} \quad & E_{\parallel 0} = E_\parallel & \\
              E_{\perp} = E_{\perp 0}\frac{B}{B_0}       \quad  & E_\parallel = 0               &
          \end{align*}
          Tokom elastičnog sudara važi zakon održanja energije:
          \begin{align*}
               & E = E'                                                                                                         \\
               & E_\perp = E_\perp' + E_\parallel' = 2E_\perp' \quad (v_\parallel' = v_\perp' \implies E_\parallel' = E_\perp') \\
               & E_\perp' = \frac{E_\perp}{2} = \frac{E_{\perp 0}}{2}\frac{B}{B_0}                                              \\
               & E_\parallel' = \frac{E_\perp}{2} = \frac{E_{\perp 0}}{2}\frac{B}{B_0}
          \end{align*}
          Nakon toga ponovo imamo vremenski promenljivo polje, pa zbog održanja magnetnog momenta ponovo dolazi do promene normalne energije (dok paralelna ostaje ista):
          \begin{align*}
              \frac{E_\perp'}{B} = \frac{E_{\perp}''}{B_0} \quad                                    & E_\parallel' = E_\parallel''                       \\
              E_\perp'' = E_\perp' \frac{B_0}{B}                            \quad                   & E_\parallel'' = \frac{E_{\perp 0}}{2}\frac{B}{B_0} \\
              E_\perp'' = \frac{E_{\perp 0}}{2} \cancel{\frac{B}{B_0}} \cancel{\frac{B_0}{B}} \quad &
          \end{align*}
          \[ E_f = E_\perp'' + E_\parallel'' = \frac{E_{\perp 0}}{2} \left(1 + \frac{B}{B_0}\right) = 5.5\si{\kilo\electronvolt} \]

          \hrulefill
    \item Dokazati da je u nehomogenom i stacionarnom magnetskom polju magnetski moment čestice konstantan.

          \underline{Rešenje}:

          Pretpostavljamo da se magnetno polje menja sporo i primarno u $z$ pravcu.
          Međutim, zbog zakona o održanju magnetskog fluksa zaključujemo da mora postojati i promena u radijalnom pravcu:
          \[ \nabla \vec{B} = 0 \implies \frac{1}{r}\frac{\partial(rB_r)}{\partial r} + \frac{1}{r}\frac{\partial B_\phi}{\partial\phi} + \frac{\partial B_z}{\partial z} = 0 \]
          Ako pretpostavimo da ne postoji azimutalna komponenta polja ($B_\phi = 0$):
          \begin{align*}
               & \frac{\partial B_z}{\partial z} = -\frac{1}{r}\frac{\partial (rB_r)}{\partial r} \\
               & rB_r = -\int r\frac{\partial B_z}{\partial z}dr
          \end{align*}
          Ako je promjena polja u $z$ pravcu dovoljno spora, u intervalu od interesa možemo smatrati da je $\frac{\partial B_z}{\partial z} \approx const$:
          \begin{align*}
               & rB_r \approx -\frac{\partial B_z}{\partial z} \int rdr = -\frac{r^2}{2} \frac{\partial B_z}{\partial z} \\
               & \implies B_r = -\frac{r}{2}\frac{\partial B_z}{\partial z}; B_z \gg B_r
          \end{align*}
          Od interesa nam je da analiziramo koji efekat nehomogeno magnetsko polje ima na naelektrisanu česticu.
          Sa poznatim njegovim vektorom sada možemo odrediti silu kojom ono deluje na česticu:
          \begin{align*}
               & \vec F = q \vec v \times \vec B = q \begin{vmatrix}
                                                         \vec e_r & \vec e_\phi & \vec e_z \\
                                                         v_r      & v_\phi      & v_z      \\
                                                         B_r      & 0           & B_z
                                                     \end{vmatrix} \\
               & F_r = q(v_\phi B_z - v_z \cdot 0)                                         \\
               & F_\phi = q(v_z B_r - v_r B_z)                                             \\
               & F_z = q(v_r \cdot 0 - B_r v_\phi)
          \end{align*}
          Već je poznato da članovi $q v_\phi B_z$ i $-qv_rB_z$ definišu Larmorovo kružno kretanje čestice oko centra vođenja.
          Analizirajmo efekat preostalih članova:
          \[ F_z = -qv_\phi B_r = \frac{qr}{2}v_\phi \frac{\partial B_z}{\partial z} \]
          Od interesa nam je da posmatramo usrednjeni efekat na sve čestice:
          \[ \left<F_z\right> = \frac{q}{2}\frac{\partial B_z}{\partial z} \left<rv_\phi\right> \]
          Član $\left<rv_\phi\right>$ je definisan Lormorovim kretanjem:
          \begin{align*}
               & \left<rv_\phi\right> = \mp R_L v_\perp \quad \text{(pozitivne čestice skreću u negativnom pravcu, i obratno)}                                                                  \\
               & \left<rv_\phi\right> = \mp \frac{v_\perp}{\omega_c}\cdot v_\perp = \mp \frac{m v_\perp^2}{|q|B} \cdot \frac{2}{2} = \mp \frac{2}{|q|}\frac{E_\perp}{B} = \mp \frac{2}{|q|} \mu \\
               & \left<F_z\right> = \frac{q}{2} \frac{\partial B_z}{\partial z} \left(\mp \frac{2}{|q|}\mu\right)                                                                               \\
               & \left<F_z\right> = -\mu \frac{\partial B_z}{\partial z} \leftarrow \text{\underline{dijamagnetska sila}}
          \end{align*}
          Do ponašanja magnetskog momenta u vremenu dolazimo polazeći od 2. Njutnovog zakona:
          \begin{align*}
               & m \frac{dv_z}{dt} = -\mu \frac{dB}{dz} = -\mu \frac{dB}{dt} \cdot \underbrace{\frac{dt}{dz}}_{\frac{1}{v_z}} \footnote{Član $\frac{dB}{dt}$ opisuje promenu magnetskog polja koje čestica \enquote{oseća}, polje samo po sebi je vremenski nepromenljivo, kao što je navedeno u postavci.} \\
               & m v_z \frac{dv_z}{dt} = -\mu \frac{dB}{dt}                                                                                                                                                                                                                                                 \\
               & \frac{d}{dt}\left(\frac{mv_z^2}{2}\right) = \frac{dE_\parallel}{dt} = -\mu \frac{dB}{dt}
          \end{align*}
          Iz zakona održanja energije dobijamo:
          \begin{align*}
               & \frac{dE}{dt} = 0                                                                      \\
               & \frac{dE_\parallel}{dt} + \frac{dE_\perp}{dt} = 0                                      \\
               & -\mu \frac{dB}{dt} + \frac{d}{dt} (\mu B) = 0                                          \\
               & -\cancel{\mu \frac{dB}{dt}} + \frac{d\mu}{dt} B + \cancel{\mu \frac{dB}{dt}} = 0       \\
               & \frac{d\mu}{dt} = 0 \implies \text{\underline{magnetski moment se ne menja u vremenu}}
          \end{align*}

          \hrulefill
    \item Plazma sa izotropnom raspodelom brzina se nalazi u magnetskom ogledalu $R_0 = 4$ (koeficijent magnetskog ogledala).
          Nema sudara, tako da čestice jednostavno beže, dok ostatak ostaje zarobljen.
          Kolika je zarobljena frakcija?

          \underline{Rešenje}:

          Iz prethodnog zadatka nam je poznato da je $\mu = \frac{E_\perp}{B} = const$ i $E = E_\perp + E_\parallel = const$.

          \noindent\hfill
          \begin{tikzpicture}
              \begin{axis}[
                      axis lines=middle,
                      xlabel={$z$},
                      xmin=-3, xmax=3,
                      ymin=-3, ymax=3,
                      xtick={-2, 2},
                      xticklabels={$a$, $b$},
                      ytick=\empty,
                      axis y line=none
                  ]
              \end{axis}
          \end{tikzpicture}
          \hfill\null\\
          Posmatramo česticu koja se kreće iz tačke $a$ sa najslabijim poljem i dolazi do tačke $b$ sa najjačim poljem.
          Da bi došlo do odbijanja, potrebno je da čestica izgubi svu paralelnu energiju dok dođe do tačke $b$, odnosno da ispunjava uslov:
          \begin{align*}
               & E_{\perp b} \ge E_b = E_a                                                      \\
               & \mu = const \implies \frac{E_{\perp a}}{B_{min}} = \frac{E_{\perp b}}{B_{max}} \\
               & E_{\perp b} = E_{\perp a} \frac{B_{max}}{B_{min}} \ge E_a                      \\
               & \frac{E_a}{E_{\perp a}} \le \frac{B_{max}}{B_{min}} = R_0
          \end{align*}
          Ovaj uslov odbijanja po energijama možemo prevesti na brzine, zato što je dato da je njihova raspodela izotropna.
          Onda možemo integracijom po prostornom uglu da zaključimo koji procenat brzina ga zadovoljava:
          \[ \frac{E_a}{E_{\perp a} = \frac{v_a^2}{v_{\perp a}^2}} \le R_0 \implies \frac{v_a}{v_{\perp a}} \le \sqrt{R_0} \]
          \hfill
          \begin{tikzpicture}
              \coordinate (O) at (0, 0);
              \coordinate (P) at (3, 0);
              \coordinate (V) at (3, 1);
              \draw[->, thick, line cap=round] O--V node[above, pos=0.5] {$\vec v$};
              \draw[->, thick, line cap=round] O--P node[below, pos=0.5] {$\vec v_\parallel$};
              \draw[->, thick, line cap=round] P--V node[right, pos=0.5] {$\vec v_\perp$};
              \pic[
                  draw,
                  angle radius=16mm,
                  angle eccentricity=0.8,
                  pic text={$\theta$}
              ] {angle = P--O--V};
          \end{tikzpicture}
          \hfill\null\\
          \begin{align*}
               & \frac{v_a}{v_{\perp a}} = \frac{v_a}{v_a \sin\theta} = \frac{1}{\sin\theta} \le \sqrt{R_0}                                  \\
               & \sin\theta \ge \frac{1}{\sqrt{R_0}}                                                                                         \\
               & \theta \ge \arcsin\left(\frac{1}{\sqrt{R_0}}\right) = \theta_{gr} \rightarrow \text{ugao koji definiše tzv. konus gubitaka}
          \end{align*}
          \noindent\hfill
          \begin{tikzpicture}[
                  >=Stealth,
              ]
              \coordinate (O) at (0, 0);
              \coordinate (P) at (2.5, 0);
              \coordinate (V) at (2.5, 1);
              \draw[->, thick, line cap=round] (-3,0)--(3,0) node[right] {$B$};
              \draw (-2.5, -1)--(2.5, 1);
              \draw (-2.5, 1)--(2.5, -1);
              \draw (-2.5, 0) ellipse (0.2cm and 1cm);
              \draw (2.5, 0) ellipse (0.2cm and 1cm);
              \pic[
                  draw,
                  angle radius=16mm,
                  angle eccentricity=0.8,
                  pic text={$\theta_{gr}$}
              ] {angle = P--O--V};
          \end{tikzpicture}
          \hfill\null\\
          Procenat čestica koji pobegne ćemo naći kao odnos prostornog ugla koji im to omogućava i ukupnog prostornog ugla:
          \begin{align*}
               & \eta_\text{loss} = \frac{\Delta\Omega}{4\pi} \cdot 2 \quad \text{(množimo sa 2 jer imamo dva konusa)}                                                                                       \\
               & d\Omega = \frac{dS}{r^2} = \frac{r^2 \sin\theta d\theta d\phi}{r^2} = \sin\theta d\theta d\phi                                                                                              \\
               & \Delta \Omega = \int d\Omega = \int_{\phi = 0}^{2\pi} \int_{\theta = 0}{\theta_{gr}} \sin\theta d\theta d\phi = 2\pi \left.(-\cos\theta)\right|_0^{\theta_{gr}} = 2\pi(1 - \cos\theta_{gr}) \\
               & \eta_\text{loss} = \frac{2\Delta\Omega}{4\pi} = 1 - \cos\theta_{gr}
          \end{align*}
          Konačno, frakciju zarobljenih čestica dobijamo kao:
          \[ \eta_\text{trap} = 1 - \eta_\text{loss} = \cos\theta_{gr} = 0.86 = 86\% \]

          \hrulefill
    \item Poznata je $z$ komponenta magnetskog polja duž ose magnetskog ogledala: \[ B_z (z) = B_0 (1 + \alpha^2 z^2) \]
          \begin{enumerate}
              \item Ako u trenutku $t = 0$ na poziciji $z = 0$ elektron ima brzinu $v^2 = 3v_\parallel^2$, pri kojoj vrednosti $z_\text{refl}$ će on biti reflektovan?
              \item Napisati jednačinu kretanja centra vođenja čestice za pravac koji je paralelan polju.
              \item Pokazati da je kretanje sinusoidalno i izračunati njegovu frekvenciju.
          \end{enumerate}

          \underline{Rešenje}:

          \begin{enumerate}
              \item Poznati su zakoni održanja:
                    \begin{align*}
                         & \mu = \frac{E_\perp}{B_z} = \frac{mv_\perp^2}{2B_z} = const \implies \frac{v_\perp^2}{B_z} = const           \\
                         & E = \frac{mv_\perp^2}{2} + \frac{mv_\parallel^2}{2} = const \implies v_\perp^2 + v_\parallel^2 = v^2 = const \\
                    \end{align*}
                    Dato je $v_{\parallel 0}$, pa prvo odredimo $v_{\perp 0}$:
                    \[ v^2 = v_{\perp 0}^2 + v_{\parallel 0}^2 = v_{\perp 0}^2 + \frac{v^2}{3} \implies v_{\perp 0}^2 = \frac{2}{3} v^2 \]
                    Sada možemo iskoristiti zakon o održanju magnetskog momenta, pri čemu imamo na umu da je u tački refleksije $v_{\parallel \text{refl}} = 0$:
                    \begin{align*}
                         & \frac{v_{\perp 0}^2}{B(0)} = \frac{v_{\perp \text{refl}^2}}{B(z_\text{refl})}                                   \\
                         & \frac{\frac{2}{3}\cancel{v^2}}{\cancel{B_0}} = \frac{\cancel{v^2}}{\cancel{B_0} (1 + \alpha^2 z_\text{refl}^2)} \\
                         & 1 + \alpha^2 z_\text{refl}^2 = \frac{3}{2}                                                                      \\
                         & \alpha^2 z_\text{refl} = \frac{1}{2}                                                                            \\
                         & z_\text{refl} = \frac{1}{\sqrt2 \alpha}
                    \end{align*}
              \item Na česticu deluje samo dijamagnetska sila $F_z$, pa je jednačina kretanja:
                    \begin{align*}
                         & m \frac{dv_z}{dt} = m \frac{d^2 z}{dt^2} = -\mu \frac{dB}{dz} = -\mu B_0 \alpha^2 2z \\
                         & \frac{d^2 z}{dt^2} + \frac{2 \mu B_0 \alpha^2}{m} z = 0
                    \end{align*}
              \item Ako uvedemo oznaku $\Omega = \sqrt{\frac{2 \mu B_0 \alpha^2}{m}}$, opšte rešenje diferencijalne jednačine možemo zapisati kao:
                    \[ z = A\sin(\Omega t) + B\cos(\Omega t) \]
                    Nepoznate konstante $A$ i $B$ nalazimo iz početnih uslova:
                    \begin{align*}
                         & z(t = 0) = 0 \implies B = 0                                                                                                                                                                                                                                \\
                         & \frac{dz}{dt}(t = 0) = v_z(t = 0) = v_{\parallel 0}                                                                                                                                                                                                        \\
                         & v_z(t) = \frac{dz}{dt} = A\Omega\cos(\Omega t)                                                                                                                                                                                                             \\
                         & v_{\parallel 0} = A\Omega                                                                                                                                                                                                                                  \\
                         & \Omega^2 = \frac{2\mu B_0 \alpha^2}{m} = \frac{2 \frac{E_{\perp 0}}{\cancel{B_0}} \cancel{B_0} \alpha^2}{m} = \frac{\cancel 2 \frac{\cancel m v_{\perp 0}^2}{\cancel{2}} \alpha^2}{\cancel m} = v_{\perp 0}^2 \alpha^2 \implies \Omega = v_{\perp 0}\alpha \\
                         & A = \frac{v_{\parallel 0}}{\Omega} = \frac{v_{\parallel 0}}{v_{\perp 0} \alpha} = \frac{\frac{\cancel{v}}{\cancel{\sqrt3}}}{\sqrt{\frac{2}{\cancel 3}} \cancel v \alpha} = \frac{1}{\sqrt 2 \alpha} = z_\text{refl}
                    \end{align*}
                    Dobili smo da je amplituda oscilovanja jednaka položaju tačke refleksije, što smo i očekivali.
          \end{enumerate}
\end{enumerate}

\end{document}